\specialchapter{术语索引}

参考编号依次表示章、节和小节（或特殊情况下表示习题）。

抽象测度、积分：IX, 3, 习题 4.

可加集合函数：IX, 3, 2.

可数可加（集合函数）：IX, 3, 2.

线性映射的伴随：IX, 附录, 2.

三角代数（上、下）：VII, 3, 3.

几乎处处：IX, 1, 9.

与测度关联的前测度：IX, 6, 1.

原子负担：IX, 3, 习题 9.

以 \(\mu\) 为基的测度：IX, 2, 2.

Bochner 定理：IX, 6, 12.

有界负担：IX, 1, 1.

有界测度：IX, 1, 2.

有界集合函数：IX, 3, 2.

Brunn--Minkowski 不等式：VII, 1, 习题 25.

强不可达基数：IX, 3, 习题 13, 脚注。

Ulam 型基数：IX, 3, 习题 11.

幺半群的特征：IX, 5, 7.

集中在子集上的负担：IX, 1, 1.

集中在子集上的测度：IX, 1, 4.

连续线性表示：VIII, 2, 1.

逆伴随（线性表示）：VIII, 2, 2.

有限测度序列的卷积积：VIII, 1, 1.

函数的卷积积：VIII, 4, 5.

可卷积函数：VIII, 4, 5.

可卷积的测度和函数：VIII, 4, 1.

可卷积、\(\varphi\)-可卷积（有限测度序列）：VIII, 1, 1.

可数可加集合函数：IX, 3, 2.

高斯测度的协方差矩阵：IX, 6, 6.

\(R^{T}\) 上的高斯前测度的协方差：IX, 6, 6.

有界测度的压缩集：VIII, 3, 习题 10.

压缩：IX, 1, 8.

\(\mathbf{GL}(n,\mathbf{K})\) 的 Iwasawa 分解：VII, 3, 3.

一个测度相对于另一个测度的密度：IX, 2, 2.

弥散负担：IX, 3, 习题 9.

测度的分解：IX, 2, 7.

基本域：VII, 2, 10.

负担：IX, 1, 1.

有界负担：IX, 1, 1.

集中在子集上的负担：IX, 1, 1.

负担的像：IX, 1, 1.

诱导负担：IX, 1, 1.

局部有界负担：IX, 1, 1.

等度连续线性表示：VIII, 2, 1.

本质上上积分：IX, 1, 2.

本质有界函数：IX, 1, 10.

本质可积函数：IX, 1, 10.

有限测度序列，\(\varphi\)-可卷积：VIII, 1, 1.

核二次型：IX, 附录, 1.

满子幺半群：IX, 5, 7.

正型函数：IX, 6, 12.

本质有界函数：IX, 1, 10.

本质可积函数：IX, 1, 10.

生成函数（某序列的）：IX, 5, 7.

可积函数：IX, 1, 10.

局部可积函数：IX, 2, 2.

局部可忽略函数：IX, 1, 4.

可测函数：IX, 1, 5.

适度函数：IX, 1, 9.

模函数（局部紧群的）：VII, 1, 3.

可忽略函数：IX, 1, 9.

普遍可测函数：IX, 2, 7.

基本域：VII, 2, 10.

G-覆盖：VII, 1, 习题 27.

G-填充：VII, 1, 习题 27.

G-铺砌：VII, 1, 习题 27.

集合 X 上的负担：IX, 3, 习题 9.

高斯测度、前测度：IX, 6, 5.

协方差矩阵为 \(C:\) 的高斯测度：IX, 6, 6.

协方差为 K 的高斯前测度：IX, 6, 6.

实希尔伯特空间上的典范高斯前测度：IX, 6, 6.

一般线性群：VII, 3, 3.

序列的生成函数：IX, 5, 7.

群，大三角（上、下）：VII, 3, 3.

群，特殊三角（上、下）：VII, 3, 3.

群，严格三角（上、下）：VII, 3, 3.

单模群：VII, 1, 3.

Haar 测度（左、右）：VII, 1, 2.

Hilbert--Schmidt 映射：IX, 附录, 2.

测度的像：IX, 2, 3.

前测度的像：IX, 6, 2.

负担的像：IX, 1, 1.

诱导负担：IX, 1, 1.

诱导测度：IX, 2, 1.

不等式，Brunn--Minkowski：VII, 1, 习题 25.

内正则集合函数：IX, 3, 2.

\(p\)-进整数：VII, 1, 6.

可积函数：IX, 1, 10.

可积集合：IX, 1, 9.

积分，抽象：IX, 3, 习题 4.

积分，本质上上：IX, 1, 2.

积分，上：IX, 1, 9.

函数的积分：IX, 1, 10.

群上的不变测度（左、右）：VII, 1, 1.

在算子群下不变的测度：VII, 1, 1.

局部紧群的逆（或射影）极限上的测度的逆（或射影）极限：VII, 1, 6.

测度的逆（或射影）极限：IX, 4, 2.

测度的逆（或射影）系统：IX, 4, 2.

等距线性表示：VIII, 2, 1.

\(\mathbf{GL}(n,\mathbf{K})\) 的 Iwasawa 分解：VII, 3, 3.

正型核：IX, 6, 6.

Lagrange 定理：VII, 1, 习题 29.

拉普拉斯变换：IX, 5, 7.

大三角群（上、下）：VII, 3, 3.

左 Haar 测度：VII, 1, 2.

局部紧群上相对不变测度的左乘子：VII, 1, 8.

群上的左不变测度：VII, 1, 1.

P. Lévy 的定理：IX, 5, 习题 13.

测度的提升：IX, 2, 4.

极限，局部紧群的逆（或射影）极限上的测度的逆（或射影）：VII, 1, 6.

极限，测度的逆（或射影）：IX, 4, 2.

线性表示，转置的、线性表示的逆伴随：VIII, 2, 2.

局部几乎处处：IX, 1, 4.

局部有界负担：IX, 1, 1.

局部有界集合函数：IX, 3, 2.

局部可积函数：IX, 2, 2.

局部可忽略函数：IX, 1, 4.

局部可忽略集合：IX, 1, 4.

映射，Hilbert--Schmidt：IX, 附录, 2.

映射，\(\mu\)-适当：IX, 2, 3.

函数空间上测度的边缘：IX, 4, 3.

总质量（前测度的）：IX, 6, 1.

轨道平均：VII, 2, 2.

可测函数：IX, 1, 5.

可测集合：IX, 1, 5.

测度：IX, 1, 2.

测度，抽象：IX, 3, 习题 4.

\(\mathbf{Q}_p\) 上的测度，归一化 Haar：VII, 1, 6.

测度，有界：IX, 1, 2.

测度，协方差矩阵为 C 的高斯：IX, 6, 6.

测度，方差为 Q 的高斯：IX, 6, 5.

测度，局部紧群上的 Haar（左、右）：VII, 1, 2.

测度，像：IX, 2, 3.

测度，诱导：IX, 2, 1.

测度，在算子群下不变（相对不变、拟不变）：VII, 1, 1.

测度，局部紧群上的左不变（右不变）：VII, 1, 1.

测度，提升的：IX, 2, 4.

测度，适度：IX, 1, 9.

测度，紧群上、离散群上的归一化 Haar：VII, 1, 3.

测度，外（集合的）：IX, 1, 9.

测度，乘积：IX, 2, 5.

测度，局部紧群上的拟不变：VII, 1, 9.

测度，局部紧群上的相对不变：VII, 1, 8.

测度，Wiener：IX, 6, 7.

测度，基为 \(\mu : \mathrm{IX}, 2, 2\) 。

测度，具有密度 \(f\)，相对于测度 \(\mu: \mathrm{IX}, 2, 2\) 。

测度空间：IX, 6, 7, 脚注 (2).

Minkowski 定理：VII, 1, 习题 27.

Minlos 定理：IX, 6, 10.

适度函数：IX, 1, 9.

适度测度：IX, 1, 9.

适度集合：IX, 1, 9.

局部紧群的模函数：VII, 1, 3.

局部紧群的模：VII, 1, 3.

自同构的模：VII, 1, 4.

局部紧群上相对不变测度的乘子（左、右）：VII, 1, 8.

乘积 \(\mathbf{G} \times \mathbf{X}\) 上的乘子；群 \(\mathbf{G}\) 与集合 \(\mathbf{X}\)，且 \(\mathbf{G}\) 作用于其上：VIII, 2, 3.

在算子群下相对不变的测度的乘子：VII, 1, 1.

可忽略函数：IX, 1, 9.

可忽略集合：IX, 1, 9.

紧群上、离散群上的归一化 Haar 测度：VII, 1, 3.

\(\mathbf{Q}_p\) 上的归一化 Haar 测度：VII, 1, 6.

核二次型：IX, 附录, 1.

核空间：IX, 6, 10.

轨道平均：VII, 2, 2.

外测度：IX, 1, 9.

\(p\)-进整数：VII, 1, 6.

正前测度、测度：IX, 1, 2.

正型函数：IX, 6, 12.

正型核：IX, 6, 6.

前测度：IX, 1, 2.

前测度，正：IX, 1, 2.

前测度，实：IX, 1, 2.

测度族的乘积：IX, 4, 3.

两个测度的乘积：IX, 2, 5.

乘积，卷积（测度与函数的）：VIII, 4, 1.

乘积，卷积（函数的）：VIII, 4, 5.

乘积，关于映射的卷积（测度的）：VIII, 1, 1.

测度的射影（或逆）极限：IX, 4, 2.

测度的射影（或逆）系统：IX, 4, 2.

Prokhorov 条件：IX, 4, 2 以及 IX, 5, 5.

前测度：IX, 6, 1.

与测度关联的前测度：IX, 6, 1.

前测度，实希尔伯特空间上的典范高斯：IX, 6, 6.

前测度，傅里叶变换：IX, 6, 3.

前测度，\(R^{T}\) 上协方差为 K 的高斯：IX, 6, 6.

前测度，方差为 Q 的高斯：IX, 6, 5.

前测度，像：IX, 6, 2.

前测度，总质量：IX, 6. 1.

\(\mu\)-适当（映射）：IX, 2, 3.

二次型，核：IX, 附录, 1.

局部紧群上的拟不变测度：VII, 1, 9.

在算子群下拟不变的测度：VII, 1, 1.

局部紧空间 X 上的测度关于在 X 上作用的群的 Haar 测度的商：VII, 2, 2.

Radon 空间：IX, 3, 3.

实测度、前测度：IX, 1, 2.

正则表示（左、右）：VIII, 2, 5.

正则化：VIII, 4, 7.

局部紧群上的相对不变测度：VII, 1, 8.

在算子群下相对不变的测度：VII, 1, 1.

表示，连续（分别连续、等度连续、等距）：VIII, 2, 1.

表示，正则（左、右）：VIII, 2, 5.

表示，酉：VIII, 2, 习题 4.

右 Haar 测度：VII, 1, 2.

局部紧群上相对不变测度的右乘子：VII, 1, 8.

群上的右不变测度：VII, 1, 1.

Sazonov 拓扑：IX, 6, 10.

分别连续（线性表示）：VIII, 2, 1.

集合函数，可加：IX, 3, 2.

集合函数，有界：IX, 3, 2.

集合函数，可数可加：IX, 3, 2.

集合函数，内正则：IX, 3, 2.

集合函数，局部有界：IX, 3, 2.

空间，核：IX, 6, 10.

空间，Radon：IX, 3, 3.

空间，强 Radon：IX, 3, 3.

特殊三角群（上、下）：VII, 3, 3.

严格三角群（上、下）：VII, 3, 3.

强不可达基数：IX, 3, 习题 13, 脚注。

强 Radon 空间：IX, 3, 3.

测度的子逆（或子射影）系统：IX, 4, 2.

子幺半群，满：IX, 5, 7.

测度族的和：IX, 1, 7.

可求和的测度族：IX, 1, 7.

测度的支集：IX, 1, 6.

系统，测度的逆（或射影）：IX, 4, 2.

系统，测度的子逆（或子射影）：IX, 4, 2.

定理，Bochner：IX, 6, 12.

定理，Lagrange：VII, 1, 习题 29.

定理，Minkowski：VII, 1, 习题 27.

定理，Minlos：IX, 6, 10.

定理，P. Lévy：IX, 5, 习题 13.

定理，Thue：VII, 1, 习题 28.

Thue 定理：VII, 1, 习题 28.

紧收敛的拓扑：IX, 5, 3.

紧拓扑：IX, 5, 3.

拓扑，Sazonov：IX, 6, 10.

拓扑，紧：IX, 5, 3.

前测度的总质量：IX, 6, 1.

一个二次型相对于另一个二次型的迹：IX, 附录, 1.

变换，傅里叶（测度的、前测度的）：IX, 6, 3.

变换，拉普拉斯（测度的）：IX, 5, 7.

转置线性表示：VIII, 2, 2.

三角代数：VII, 3, 3.

三角群（大、严格、特殊）：VII, 3, 3.

Ulam 基数：IX, 3, 习题 11.

Ulam 超滤子：IX, 3, 习题 11.

单模群：VII, 1, 3.

酉表示：VIII, 2, 习题 4.

普遍可测函数：IX, 2, 7.

普遍可测集合：IX, 3, 3.

上积分：IX, 1, 9.

测度的方差：IX, 6, 5.

Wiener 测度：IX, 6, 7.

\specialchapter{第七章的主要公式}

\specsection{\texorpdfstring{有关公式：\(\pmb{\gamma}(s)\) 和 \(\delta(s)\)}{有关公式：\textbackslash pmb\{\textbackslash gamma\}(s) 和 \textbackslash delta(s)}}

设 \(G\) 是一个拓扑群，它在局部紧空间 \(X\) 上通过 \((s, x) \mapsto sx\) 连续地左作用。

\(\begin{array}{ll}\pmb {\gamma}(s)x = sx & (s\in \mathrm{G},x\in \mathrm{X})\\ \pmb {\gamma}(st) = \pmb {\gamma}(s)\pmb {\gamma}(t) & (s,t\text{ 属于 G})\\ (\pmb {\gamma}(s)f)(x) = f(s^{-1}x) & (f\text{ X 上的函数})\\ \langle f,\pmb {\gamma}(s)\mu \rangle = \langle \pmb {\gamma}(s^{-1})f,\mu \rangle & (\mu \text{ X 上的测度})\\ d(\pmb {\gamma}(s)\mu)(x) = d\mu (s^{-1}x)\\ (\pmb {\gamma}(s)\mu)(\mathrm{A}) = \mu (s^{-1}\mathrm{A}) & (\mathrm{A}\text{ 是一个 }\pmb {\gamma}(s)\mu \text{-可积集}) \end{array}\)

若 \(\mu\) 是乘子为 \(\chi\) 的相对不变测度，

\[
\begin{array}{l} \pmb {\gamma} (s) \mu = \chi (s) ^ {- 1} \mu \\ d \mu (s x) = \chi (s) d \mu (x). \end{array}
\]

设 \(G\) 是一个拓扑群，它在局部紧空间 \(X\) 上通过 \((s, x) \mapsto xs\) 连续地右作用。

\(\pmb {\delta}(s)x = xs^{-1}\) \(\pmb {\delta}(st) = \pmb {\delta}(s)\pmb {\delta}(t)\) \((\pmb {\delta}(s)f)(x) = f(xs)\) \(\langle f,\pmb {\delta}(s)\mu \rangle = \langle \pmb {\delta}(s^{-1})f,\mu \rangle\) \(d(\pmb {\delta}(s)\mu)(x) = d\mu (xs)\) \((\pmb {\delta}(s)\mu)(\mathrm{A}) = \mu (\mathrm{As}).\)

若 \(\mu\) 是乘子为 \(\chi'\) 的相对不变测度，

\[
\begin{array}{l} \pmb {\delta} (s) \mu = \chi' (s) \mu \\ d \mu (x s) = \chi' (s) d \mu (x). \end{array}
\]

\specsection{关于 Haar 测度的公式}

设 \(G\) 是一个局部紧群，\(\Delta\) 是它的模，\(\mu\) 是左 Haar 测度，\(\nu\) 是右 Haar 测度。

\begin{enumerate}
\def\labelenumi{\arabic{enumi})}
\tightlist
\item
  有
\end{enumerate}

\[
\begin{array}{l l l l} \boldsymbol {\gamma} (s) \mu = \mu & \quad \boldsymbol {\delta} (s) \mu = \Delta (s) \mu & \quad \check {\mu} = \Delta^ {- 1} \cdot \mu \\ d \mu (s x) = d \mu (x) & \quad d \mu (x s) = \Delta (s)   d \mu (x) & \quad d \mu (x ^ {- 1}) = \Delta (x) ^ {- 1}   d \mu (x). \end{array}
\]

若 \(f\) 是 \(\mu\)-可积的，

\[
\begin{array}{c} \int f (s x) d \mu (x) = \int f (x) d \mu (x) \\ \int f (x s) d \mu (x) = \Delta (s) ^ {- 1} \int f (x) d \mu (x) \\ \int f (x ^ {- 1}) \Delta (x) ^ {- 1} d \mu (x) = \int f (x) d \mu (x). \end{array}
\]

若 \(\mathrm{A}\subset \mathrm{G}\) 是 \(\mu\)-可积的，

\[
\mu (s \mathrm{A}) = \mu (\mathrm{A}) \quad \mu (\mathrm{A} s) = \Delta (s) \mu (\mathrm{A}).
\]

\begin{enumerate}
\def\labelenumi{\arabic{enumi})}
\setcounter{enumi}{1}
\tightlist
\item
  有
\end{enumerate}

\[
\begin{array}{l l l l} \boldsymbol {\delta} (s) \nu = \nu & \boldsymbol {\gamma} (s) \nu = \Delta (s) \nu & \stackrel {{\vee}} {{\nu}} = \Delta \cdot \nu \\ d \nu (x s) = d \nu (x) & d \nu (s ^ {- 1} x) = \Delta (s)   d \nu (x) & d \nu (x ^ {- 1}) = \Delta (x)   d \nu (x). \end{array}
\]

若 \(f\) 是 \(\nu\)-可积的，

\[
\begin{array}{c} \int f (x s) d \nu (x) = \int f (x) d \mu (x) \qquad \int f (s x) d \nu (x) = \Delta (s) \int f (x) d \nu (x) \\ \int f (x ^ {- 1}) \Delta (x) d \nu (x) = \int f (x) d \nu (x). \end{array}
\]

若 \(\mathbf{A}\subset \mathbf{G}\) 是 \(\nu\)-可积的，

\[
\nu (\mathrm{A} s) = \nu (\mathrm{A}) \quad \nu (s \mathrm{A}) = \Delta (s ^ {- 1}) \nu (\mathrm{A}).
\]

\begin{enumerate}
\def\labelenumi{\arabic{enumi})}
\setcounter{enumi}{2}
\tightlist
\item
  \(\nu\) 与 \(\Delta^{-1} \cdot \mu\) 成比例，\(\mu\) 与 \(\Delta \cdot \nu\) 成比例。
\end{enumerate}

\specialchapter{卷积积存在的充分条件}

I. --- 两个测度的卷积积 \(\mu * \nu\) 存在的情形：

\begin{enumerate}
\def\labelenumi{(\alph{enumi})}
\tightlist
\item
  * 由连续映射 \(\varphi : \mathbf{X} \times \mathbf{Y} \to \mathbf{Z}\) 定义：
\end{enumerate}

\(\mu, \nu\) 有界（于是 \(\mu * \nu\) 有界，且 \(\| \mu * \nu \| \leqslant \| \mu \| \cdot \| \nu \|\)）。

\(\mu, \nu\) 具有紧支集（于是 \(\mu * \nu\) 具有紧支集，且 \(\text{Supp}(\mu * \nu) \subset \varphi(\text{Supp } \mu \times \text{Supp } \nu)\)）。

\begin{enumerate}
\def\labelenumi{(\alph{enumi})}
\setcounter{enumi}{1}
\item
  * 由一个在空间上连续左作用的群定义：\(\mu\) 具有紧支集，\(\nu\) 任意。
\item
  * 由群 G 中的乘法定义：
\end{enumerate}

两个测度中至少有一个具有紧支集。

\(\mu, \nu\) 属于 \(\mathcal{M}^{\rho}(\mathrm{G})\)（于是 \(\mu * \nu \in \mathcal{M}^{\rho}(\mathrm{G})\)，且 \(\| \mu * \nu \|_{\rho} \leqslant \| \mu \|_{\rho} \| \nu \|_{\rho}\)）。

\begin{enumerate}
\def\labelenumi{\Roman{enumi}.}
\setcounter{enumi}{1}
\tightlist
\item
  --- 测度与函数的卷积积 \(\mu*f\) 存在的情形：
\end{enumerate}

\begin{enumerate}
\def\labelenumi{(\alph{enumi})}
\tightlist
\item
  * 由群 G 在赋有测度 \(\beta \geqslant 0\) 的空间 X 上连续左作用定义，并满足 \(\pmb{\gamma}(s)\beta = \chi(s^{-1},\cdot)\beta\)，其中 \(\chi\) 连续：
\end{enumerate}

\(\mu\) 具有紧支集，f 局部 \(\beta\)-可积（若 f 连续，则 \(\mu * f\) 连续；若 f 连续且具有紧支集，则 \(\mu * f\) 连续且具有紧支集）。

G 在 X 中真作用，\(f \in \mathcal{K}(X)\)（\(\mu * f\) 连续）。

\begin{enumerate}
\def\labelenumi{(\alph{enumi})}
\setcounter{enumi}{1}
\tightlist
\item
  \(\chi(s, \cdot)\) 有界；令 \(\rho(s) = \sup_{x \in X} \chi(s^{-1}, x)\)：
\end{enumerate}

\(\mu \in \mathcal{M}^{\rho}(\mathrm{G})\)，\(f\in \mathrm{L}^{\infty}(\mathrm{X},\beta)\)（于是 \(\mu *f\in \mathrm{L}^{\infty}(\mathrm{X},\beta)\)；若 \(f\in \mathcal{C}^{\infty}(\mathrm{X})\)，则 \(\mu *f\in \mathcal{C}^{\infty}(\mathrm{X})\)；若 \(f\in \overline{\mathcal{K}(\mathrm{X})}\)，则 \(\mu *f\in \overline{\mathcal{K}(\mathrm{X})}\)）。

\(\mu \in \mathcal{M}^{\rho^{1 / q}}(\mathrm{G})\)，\(f\in \mathrm{L}^p (\mathrm{X},\beta)\)，其中 \(1 / p + 1 / q = 1\)（于是 \(\mu *f\in \mathrm{L}^p (\mathrm{X},\beta)\)，且 \(\| \mu *\boldsymbol {f}\| _p\leqslant \| \mu \|_{\rho^{1 / q}}\| f\| _p\)）。

\begin{enumerate}
\def\labelenumi{\Roman{enumi}.}
\setcounter{enumi}{2}
\tightlist
\item
  --- 卷积积 \(f * g\) 存在的情形（两个函数关于 \(\beta\) 局部可积；\(\beta\) 是满足 \(\geqslant 0\) 的、定义在群 \(G\) 上的相对不变测度，左、右乘子分别为 \(\chi\) 和 \(\chi'\)）：
\end{enumerate}

\(f\) 或 \(g\) 连续，\(f\) 或 \(g\) 具有紧支集（于是 \(f*g\) 连续；若 \(f,g\) 属于 \(\mathcal{H}(\mathrm{G})\)，则 \(f*g\in\mathcal{H}(\mathrm{G})\)）。

\(f\chi^{-1/q} \in L^{1}(G,\beta)\) 且 \(g \in L^{p}(G,\beta)\)，其中 \(1/p + 1/q = 1\)（于是 \(f*g \in L^{p}(G,\beta)\)，且 \(\|f*g\|_{p} \leqslant \|f\chi^{-1/q}\|_{1} \|g\|_{p}\)）。\(f \in L^{p}(G,\beta)\) 且 \(g\chi'^{-1/q} \in L^{1}(G,\beta)\)（于是 \(f*g \in L^{p}(G,\beta)\)，且 \(\|f*g\|_{p} \leqslant \|f\|_{p} \|g\chi'^{-1/q}\|_{1}\)）。\(f\chi^{-1} \in L^{1}(G,\beta)\) 且 \(g \in \mathcal{C}^{\infty}(G)\)（相应地为 \(\overline{\mathcal{K}(G)}\)）（于是 \(f*g \in \mathcal{C}^{\infty}(G)\)（相应地为 \(\overline{\mathcal{K}(G)}\)））。\(f \in \mathcal{C}^{\infty}(G,\beta)\)（相应地为 \(\overline{\mathcal{K}(G)}\)）且 \(g\chi'^{-1} \in L^{1}(G,\beta)\)（于是 \(f*g \in \mathcal{C}^{\infty}(G)\)（相应地为 \(\overline{\mathcal{K}(G)}\)））。\(f \in L^{p}(G,\beta), g \in L^{q}(G, \overset{\vee}{\beta})\)，其中 \(1/p + 1/q = 1, 1 < p < +\infty, \beta\) 左不变（于是 \(f*g \in \overline{\mathcal{K}(G)}\)，且 \(\|f*g\|_{\infty} \leqslant \|f\|_{p} \| \overset{\vee}{g}\|_{q}\)）。
